% LLM Policy Context Summary — LaTeX for research article
%
% Preamble (add to your document):
%   \usepackage{booktabs}
%   \usepackage{tikz}
%   \usetikzlibrary{arrows.meta, positioning}
%
% Or use as standalone: place the section content inside your article.
%
% ----------------------------------------------------------------------------

\section{LLM Policy Context and Data Flow}
\label{sec:llm-context}

Figure~\ref{fig:llm-dataflow} shows the data flow for the StellCoilBench LLM-based case proposer.
The system prompt is assembled from the surface catalog, domain context document, and base instructions.
The user prompt is built from the context payload produced by \texttt{build\_context}, which aggregates successful runs from submissions, failures from \texttt{autopilot\_failures.json}, and curated baseline cases.

\begin{figure}[t]
  \centering
  \begin{tikzpicture}[
    node distance=0.5em,
    box/.style={rectangle, draw, rounded corners=2pt, minimum height=1.8em, font=\small\tt},
    arr/.style={-{Stealth[length=2pt]}}
    ]
    \node[box, align=center] (src) {Data sources\\(policy, catalog, llm\_context, submissions, failures, cases)};
    \node[box, right=3em of src] (build) {build\_context};
    \node[box, align=center, right=3em of build] (ctx) {Context payload\\(run\_cards, postmortems, margins, trend)};
    \node[box, align=center, below=2.5em of ctx, xshift=-4em] (sys) {System prompt\\(instructions + catalog + domain doc)};
    \node[box, align=center, right=0.5em of sys] (usr) {User prompt\\(from context)};
    \node[box, right=3em of usr] (llm) {LLM};
    \node[box, right=2em of llm] (act) {JSON actions};
    \node[box, right=2em of act] (out) {CI cases};

    \draw[arr] (src) -- (build);
    \draw[arr] (build) -- (ctx);
    \draw[arr] (ctx.south) -- ++(0,-0.8em) -| (usr.north);
    \draw[arr] (sys) -- (llm);
    \draw[arr] (usr) -- (llm);
    \draw[arr] (llm) -- (act);
    \draw[arr] (act) -- (out);

    \node[font=\scriptsize, gray, above=0.2em of sys] {system};
    \node[font=\scriptsize, gray, above=0.2em of usr] {user};
  \end{tikzpicture}
  \caption{Data flow for the LLM case proposer. Data sources feed \texttt{build\_context}, which produces a context payload (run cards, postmortems, margins, score trend). The system prompt combines base instructions with the surface catalog and domain document; the user prompt is built from the context. The LLM returns JSON actions, which are converted to validated CI case configurations.}
  \label{fig:llm-dataflow}
\end{figure}

\subsection{System Prompt Components}

The system prompt is constructed from three components (Table~\ref{tab:system-prompt}).

\begin{table}[t]
  \centering
  \caption{System prompt components}
  \label{tab:system-prompt}
  \begin{tabular}{@{}llp{6.5cm}@{}}
    \toprule
    \textbf{Source} & \textbf{File} & \textbf{Content} \\
    \midrule
    Base instructions & \texttt{llm\_endpoints} & Role, JSON schema (mutate/explore), rules, strategy guidance \\
    Surface catalog & \texttt{surface\_catalog.json} & Per-surface: file, short\_name, minor\_radius\_m, a0, nfp, stellsym, typical\_ncoils, typical\_orders, notes \\
    Domain context & \texttt{llm\_context.md} & Optimization guide, objective terms, failure modes, metrics, threshold scaling, literature excerpts (7 refs) \\
    \bottomrule
  \end{tabular}
\end{table}

\subsection{User Prompt Sections}

The user prompt is assembled from the context payload (Table~\ref{tab:user-prompt}).

\begin{table}[t]
  \centering
  \caption{User prompt sections (from \texttt{build\_context} output)}
  \label{tab:user-prompt}
  \begin{tabular}{@{}llcl@{}}
    \toprule
    \textbf{Section} & \textbf{Context key} & \textbf{Max} & \textbf{Description} \\
    \midrule
    Request & --- & --- & ``Propose $N$ cases for the next optimization batch'' \\
    Prior reasoning & \texttt{prior\_reasoning} & 50 & Past proposal reasoning blocks \\
    Policy constraints & \texttt{policy.exploration} & --- & Allowed surfaces, ncoils, order, threshold ranges \\
    Top parent runs & \texttt{run\_cards} & 10 & Formatted run cards from \texttt{make\_run\_card} \\
    Failure stats & \texttt{failure\_stats} & --- & fail\_rate, failure\_classes \\
    Postmortems & \texttt{postmortems} & 5 & Failure analyses from \texttt{make\_postmortem} \\
    Surface exploration & \texttt{surface\_exploration\_counts} & --- & Runs per surface \\
    Constraint margins & \texttt{margin\_summary} & --- & Violation rate and median margin per constraint \\
    Score trend & \texttt{score\_trend} & --- & improving/plateaued/regressing, best score, best feasible \\
    Baseline cases & \texttt{baseline\_cases} & 3/surface & Curated configs from \texttt{cases/*.yaml} \\
    \bottomrule
  \end{tabular}
\end{table}

\subsection{Context Payload Structure}

Table~\ref{tab:context-payload} enumerates the keys produced by \texttt{build\_context}.

\begin{table}[t]
  \centering
  \caption{\texttt{build\_context} output structure}
  \label{tab:context-payload}
  \begin{tabular}{@{}llp{6cm}@{}}
    \toprule
    \textbf{Key} & \textbf{Type} & \textbf{Description} \\
    \midrule
    \texttt{policy} & dict & batch\_size, exploit\_fraction, resource\_caps \\
    \texttt{failure\_stats} & dict & fail\_rate, failure\_reasons, failure\_classes \\
    \texttt{top\_parents} & list & Top feasible runs by total\_score \\
    \texttt{recent\_config\_hashes} & list & Hashes of last $\sim$50 configs for novelty check \\
    \texttt{surface\_exploration\_counts} & dict & Surface name $\to$ run count \\
    \texttt{run\_cards} & list & Formatted strings from \texttt{make\_run\_card} (top 10) \\
    \texttt{postmortems} & list & Formatted strings from \texttt{make\_postmortem} (up to 5) \\
    \texttt{total\_completed} & int & Total number of summaries \\
    \texttt{margin\_summary} & dict & Per constraint: violation\_rate, median\_margin, $n$ \\
    \texttt{score\_trend} & dict & best\_score, best\_feasible\_score, trend, feasible\_count \\
    \texttt{baseline\_cases} & list & Formatted strings from \texttt{cases/*.yaml} \\
    \bottomrule
  \end{tabular}
\end{table}

\subsection{Run Card and Postmortem Formats}

Run cards (Table~\ref{tab:run-card}) summarize each successful run for the LLM. Postmortems (Table~\ref{tab:postmortem}) provide structured failure analyses.

\begin{table}[t]
  \centering
  \caption{Run card fields (\texttt{make\_run\_card})}
  \label{tab:run-card}
  \begin{tabular}{@{}lll@{}}
    \toprule
    \textbf{Field} & \textbf{Source} & \textbf{Example} \\
    \midrule
    case\_id & summary & 02-10-2026\_14-30 \\
    status & success & SUCCESS / FAILED \\
    score, iterations, walltime & summary & 1.2e-6\,$|$\,450\,$|$\,120s \\
    surface, ncoils, order & case\_config & input.QA\,$|$\,4\,$|$\,8 \\
    thresholds & coil\_objective\_terms & cc=1.2\,m, cs=2.0\,m \\
    metrics & summary.metrics & CC separation, curvature, MSC, force \\
    margins & summary.margins & Tight margins, violated constraints \\
    failure info & summary & failure\_class --- reason \\
    \bottomrule
  \end{tabular}
\end{table}

\begin{table}[t]
  \centering
  \caption{Postmortem fields (\texttt{make\_postmortem})}
  \label{tab:postmortem}
  \begin{tabular}{@{}p{4cm}p{7cm}@{}}
    \toprule
    \textbf{Field} & \textbf{Description} \\
    \midrule
    failure\_class & min\_sep\_violation, line\_search\_fail, timeout, nan\_in\_objective, vmec\_nonconverged, validation \\
    failure\_reason & Truncated reason string \\
    suggestions & Rule-based suggestions per failure\_class \\
    negative\_margins & Constraint violations with values \\
    \bottomrule
  \end{tabular}
\end{table}

\subsection{LLM Action Schema and Policy Parameters}

The LLM returns a JSON array of actions (Table~\ref{tab:actions}). Policy parameters (Table~\ref{tab:policy}) control batch size, context limits, and exploration bounds.

\begin{table}[t]
  \centering
  \caption{LLM action types and output schema}
  \label{tab:actions}
  \begin{tabular}{@{}llp{5.5cm}@{}}
    \toprule
    \textbf{Action} & \textbf{Required fields} & \textbf{Optional overrides} \\
    \midrule
    mutate & parent\_id, reasoning & ncoils, order, surface, cc/cs/curvature/msc/length/force/torque\_threshold \\
    explore & surface, ncoils, order, reasoning & thresholds dict (same keys as overrides) \\
    \bottomrule
  \end{tabular}
\end{table}

\begin{table}[t]
  \centering
  \caption{Policy parameters (\texttt{proposer\_policy.yaml})}
  \label{tab:policy}
  \begin{tabular}{@{}lll@{}}
    \toprule
    \textbf{Section} & \textbf{Key} & \textbf{Default} \\
    \midrule
    batch & batch\_size & 10 \\
    batch & exploit\_fraction & 0.5 \\
    batch & top\_k\_parents & 10 \\
    \addlinespace
    resource\_caps & max\_total\_iterations, max\_maxiter & 5000 \\
    resource\_caps & timeout\_minutes\_min/max/default & 1, 60, 30 \\
    \addlinespace
    llm\_proposer & max\_run\_cards & 10 \\
    llm\_proposer & max\_postmortems & 5 \\
    llm\_proposer & max\_prior\_reasoning\_batches & 50 \\
    llm\_proposer & temperature & 0.3 \\
    \addlinespace
    fourier\_continuation & enabled, orders & true, [4, 8, 16] \\
    \addlinespace
    exploration & surfaces & From policy (e.g.\ muse.focus) \\
    exploration & ncoils\_choices, order\_choices & [3, 4, 5, 6], [4] \\
    exploration & threshold ranges & length, cc, cs, curvature, msc, force \\
    \addlinespace
    guardrails & sliding\_window & 30 \\
    guardrails & max\_fail\_rate & 0.6 \\
    guardrails & max\_common\_failure\_count & 12 \\
    guardrails & max\_critical\_class\_count & 10 \\
    \addlinespace
    safe\_mode & threshold & 0.35 \\
    safe\_mode & max\_iterations\_cap & 5000 \\
    \addlinespace
    cooldown & batches\_after\_guardrail & 3 \\
    cooldown & write\_pause\_file & true \\
    \bottomrule
  \end{tabular}
\end{table}

\subsection{Inputs and Outputs}

\textbf{Inputs:} \texttt{policy/proposer\_policy.yaml}, \texttt{knowledge/surface\_catalog.json}, \texttt{knowledge/llm\_context.md}, \texttt{submissions/<surface>/auto/<case\_id>/results.json} and \texttt{case.yaml}, \texttt{policy/autopilot\_failures.json}, \texttt{cases/*.yaml}.

\textbf{Outputs:} A JSON array of actions \texttt{[\{type, parent\_id$|$surface, overrides$|$thresholds, reasoning\}, ...]} converted to CI case dicts via \texttt{apply\_llm\_action} and validated with \texttt{validate\_ci\_case}.
